\section{研究課題と仮説}\label{研究課題と仮説}

先行研究から，本研究では音楽嗜好の共有を通じた自己開示の場としてカラオケ環境を利用する．特に，自己開示の深さを考慮して，段階的に自己開示を促すシステムの導入により，コミュニケーションを阻害することなく効果的な自己開示が促進される可能性がある．この知見に基づき，以下の研究課題と，その課題を明らかにするための仮説を設定する．

\subsubsection{課題:カラオケ空間における音楽嗜好の段階的な自己開示が受容性に及ぼす影響}

WeberやReisらの研究では，相手の経験や感情を理解し受け入れる受容性が，良好な対人関係の構築に不可欠であることが示唆されているが，音楽嗜好の共有を通じた自己開示が受容に与える影響については十分な検討がされていない．そこで本研究では，カラオケ空間における段階的な自己開示が他者受容性に及ぼす影響を実証的に検証する．

\subsubsection{仮説1:低いレベルから高いレベルへと段階的に自己開示を行うことで，深いレベルの自己開示に対する受容度が高くなる}

仮説1は，社会的浸透理論に基づいており，対人関係の発展過程において自己開示が段階的に深まることを前提としている．突発的な深い自己開示は相手が受容できない可能性があるが，浅い自己開示から始めることにより，被開示者は開示者の個人的な経験や感情をより自然に受容できると考えられる．

\subsubsection{仮説2:音楽への認知度が低い，または興味度が低い場合でも，自己開示の受容度は大きく低下しない}

音楽嗜好の共有を通じた自己開示において，被開示者が必ずしも同じ楽曲への認知や興味を持っているとは限らない．しかし，段階的な自己開示手法を用いることで，音楽そのものへの興味の有無に関わらず，開示者の経験や感情を理解し受容できると考えられる．

\subsubsection{仮説3:セッションを終えた2者は，段階的に自己開示を促していない方に比べて相手に対する受容性が向上する}

適切な段階を経た自己開示は，開示者と被開示者の双方に安心感をもたらし，相手の経験や感情に対する理解と受容を深める．これにより，セッション後の2者間における相互の受容性が，段階的な支援がない場合と比較して向上すると考えられる．

\vskip\baselineskip

これらの仮説を検証するためには，段階的な手法の効果を明確に評価できる環境から始めることが重要である．カラオケ環境では，歌唱による本人の緊張感や，歌唱の技術レベルといった個人特性が大きく影響する．また，歌唱そのものの楽しさや，共同で場を作り上げることによる自然な親密度の向上など，多くの要因が複雑に絡み合っている．これらの要因は自己開示の促進効果を正確に測定することを困難にする可能性がある．そこで本研究ではまず，カラオケの環境に近しい状態で歌唱という要素を除いた「音楽共有セッション」を行うこととする．音楽共有セッションにおいて自己開示を段階的に促進する楽曲推薦システムと，段階的な促進を行わないシステムを用いた比較実験を行う．これにより，音楽を介したコミュニケーションにおける自己開示と受容性の基礎的な知見を明らかにした上で，将来的に歌唱を伴うカラオケ環境特有の要因を考慮した検証へと発展させることを目指す．